\section{Security Reduction: PoML to the Bitcoin Backbone} 
\label{sec:reduction}

In this section, we formally show that the PoML consensus mechanism inherits the fundamental security properties established by Garay, Kiayias, and Leonardos~\cite{GKL20} for the Bitcoin backbone protocol---namely \emph{common prefix}, \emph{chain quality}, and \emph{chain growth}.

The central idea is to show that the \emph{composite operation} of ML inference, proof generation, and hash evaluation plays the same structural role in PoML that the random oracle~$H$ plays in Bitcoin. In Bitcoin, the only way to obtain a lottery ticket (a potential valid block) is to query the random oracle; no shortcut exists. In PoML, we establish an analogous property: the only way to obtain a lottery ticket is to execute the full pipeline of inference, proof, and hash evaluation---effectively querying an \emph{inference oracle}. Once this correspondence is in place, the backbone's $q$-bounded random oracle framework applies directly, and all three backbone guarantees carry over with suitably adjusted parameters.

We adopt the notation of~\cite{GKL20} throughout (cf.\ Table~1 therein). In particular:
\begin{itemize}
    \item $\kappa$: the security parameter and length of the hash function output (i.e., $H : \{0,1\}^* \to \{0,1\}^\kappa$);
    \item $\lambda$: the tail-bounds parameter, i.e., the minimum number of consecutive rounds over which concentration of the random variables $X$, $Y$, $Z$ is required (cf.\ Definition~9 in~\cite{GKL20}), with $\lambda \geq 2/f$;
    \item $n$: total number of parties (miners), of which $t$ are controlled by the adversary;
    \item $\gamma$: honest majority slack, $t/(n-t) \leq 1 - \gamma$ (denoted $\delta$ in~\cite{GKL20}; we reserve $\delta$ for the amortizability parameter of Definition~\ref{def:delta-reuse});
    \item $f$: probability that at least one honest party succeeds in finding a PoW in a round;
    \item $\varepsilon \in (0,1)$: quality of concentration of random variables in typical executions;
    \item $L$: total runtime of the system (in rounds);
    \item $Q = qnL$: total number of random oracle queries;
    \item $\mathcal{C}^{\lceil k}$: the chain obtained by removing the last $k$ blocks from $\mathcal{C}$; $\mathcal{C}_1 \preceq \mathcal{C}_2$ denotes that $\mathcal{C}_1$ is a prefix of $\mathcal{C}_2$.
\end{itemize}


\subsection{Reduction Framework}

Recall that the backbone protocol operates in rounds, with $n$ parties of which at most $t$ are corrupted. A \textit{block} is defined in the form $B = (s,x,ctr)$, where $s$ is the previous block hash, $x$ is the block content, and $ctr$ is a counter of queries to the random oracle $H$. $B$ should satisfy the $\mathsf{validblock}^T_q$ predicate, defined as
\[
(H(ctr,G(s,x))<T) \wedge (ctr \leq q).
\]
Crucially, each flat-model party (i.e. party with same amount of computational resource assumed) is limited to $q$ queries to the random oracle per round, and the only way to produce a valid block is to query this oracle---there is no shortcut.

In PoML, the role of the random oracle is replaced by an \emph{inference oracle}: each completed ML inference coupled with a valid ZKP and a hash evaluation constitutes one query to this oracle, yielding one ``lottery ticket.''  We denote by $q_{\mathsf{ML}}$ the maximum number of such oracle queries a single flat-model honest party can complete per round, and by $D$ the PoML difficulty parameter (the lottery threshold, i.e., a block candidate is valid if $H_i < D$). We define a block in PoML as $B = (s, x, \mathcal{X}, \mathcal{Y}, \Pi, \Pi^{(2)})$.  Our $\mathsf{validblock}^D_{q_{\mathsf{ML}}}$ predicate is defined as:
\[
\begin{aligned}
(& H(G(s,x), \Pi) < D) \\
\wedge\;& (1 \le |\mathcal{X}| = |\mathcal{Y}| = |\Pi| = |\Pi^{(2)}| \le q_{\mathsf{ML}}) \\
\wedge\;& (\forall i,\ 
\mathsf{Verify}(\mathsf{CRS},\mathsf{stmt}_i^{(1)},\pi_i^{(1)})=1 ) \\
\wedge\;& (\forall i,\ 
\mathsf{Verify}(\mathsf{CRS},\mathsf{stmt}_i^{(2)},\pi_i^{(2)})=1 ).
\end{aligned}
\]


We now show that this mechanism can be mapped into the backbone's $q$-bounded framework---with the inference oracle taking the place of the random oracle---and thus inherits all its properties.

Our reduction proceeds in five steps:
\begin{enumerate}[nosep]
    \item Bound the number of inference-proof pairs per round to $q_{\mathsf{ML}}$ via amortisation resistance.
    \item Establish a bijection between proof tuples and lottery evaluations.
    \item Establish that the PoML lottery is a $q_{\mathsf{ML}}$-bounded \emph{inference oracle}.
    \item Derive the success probability $f$ for PoML.
    \item Verify typicality of executions.
\end{enumerate}

\subsection{Step 1: Bounding Inference-Proof Pairs}
\label{subsec:step0}

In Bitcoin's backbone protocol, security relies on the assumption that each flat-model party can make at most $q$ queries per round to the random oracle $H$. This bound assumes the uniform cost of each hash evaluation.

In PoML, the analogous bound must hold for \emph{inference-proof pairs}: each pair consists of (i) a model inference $M_\theta(c,r)$ and (ii) the corresponding ZKP $\pi^{(1)}$. Recall from Section~\ref{subsec:model-class} that the model class $\mathcal{K}$ consists of models with \emph{uniform query cost} (Definition~\ref{def:uniform-cost}) and \emph{$\delta$-amortizability} (Definition~\ref{def:delta-reuse}). We show that both the inference and proof generation components inherit $\delta$-amortizability, bounding the number of inference-proof pairs per round to $q_\mathsf{ML}$.

\begin{theorem}[$q_{\mathsf{ML}}$-bounded inference-proof pairs]
\label{thm:amort-resist}
Let $M_\theta \in \mathcal{K}$ be a PoML-compatible model (Definition~\ref{def:model}) with $\delta$-amortizability parameter satisfying
\begin{equation}\label{eq:delta-bound}
    \delta \;<\; \frac{1}{q_{\mathsf{ML}} + 1},
\end{equation}
and let the proof system be a non-interactive zk-SNARK. Then, in every round, each flat-model party can complete at most $q_{\mathsf{ML}}$ inference-proof pairs, except with negligible probability in $\kappa$.
\end{theorem}

\begin{proof}
Let $T_M^{\mathsf{inf}}$ and $T_M^{\mathsf{pf}}$ denote the cost of a single model inference and a single proof generation, respectively, and let $T_M = T_M^{\mathsf{inf}} + T_M^{\mathsf{pf}}$ be the cost per inference-proof pair. Let $T_{\mathsf{round}}$ denote the computation budget of a flat-model party per round, so the honest count is $q_{\mathsf{ML}} = \lfloor T_{\mathsf{round}} / T_M \rfloor$ pairs.

\paragraph{Inference cost.}
 By Definition~\ref{def:delta-reuse}, any PPT adversary requires at least $(1-\delta) \cdot T_M^{\mathsf{inf}}$ cost per inference, regardless of preprocessing or cached state.

\paragraph{Proof generation cost.}
In a non-interactive zk-SNARK, the prover's computation is a deterministic function of the common reference string (CRS), the witness, and the Fiat--Shamir challenges derived from the transcript. The witness for pair~$i$ consists of all intermediate wire values of $M_\theta(c_i, r_i)$, which by the $\delta$-amortizability of~$\mathcal{K}$ differ from those of any prior pair in at least a $(1-\delta)$-fraction of entries, with overwhelming probability. Since the prover's per-constraint computation depends on the witness values at that constraint, at most a $\delta$-fraction of the prover's work can be reused across pairs. Therefore proof generation is also $\delta$-amortizable: any adversary requires at least $(1-\delta) \cdot T_M^{\mathsf{pf}}$ per proof.

\paragraph{Combined bound.}
The adversary's per-pair cost is at least $(1-\delta) \cdot T_M^{\mathsf{inf}} + (1-\delta) \cdot T_M^{\mathsf{pf}} = (1-\delta) \cdot T_M$, yielding at most
\[
    \left\lfloor \frac{T_{\mathsf{round}}}{(1-\delta) \cdot T_M} \right\rfloor 
    = \left\lfloor \frac{q_{\mathsf{ML}}}{1 - \delta} \right\rfloor
\]
pairs. Under condition~\eqref{eq:delta-bound}, $q_{\mathsf{ML}} / (1-\delta) < q_{\mathsf{ML}} + 1$, so the floor equals $q_{\mathsf{ML}}$ and the adversary gains no additional inference-proof pair.
\end{proof}

In practical PoML deployments where $q_{\mathsf{ML}}$ is small, the bound~\eqref{eq:delta-bound} on $\delta$ is not prohibitively tight. That stochastic denoising models satisfy $\delta$-amortizability with $\delta = \mathsf{negl}(\kappa)$ is established in Appendix~\ref{app:cia}.


\subsection{Step 2: Proof-to-Lottery Bijection}

In Step~1, we established that each flat-model party can produce at most $q_{\mathsf{ML}}$ inference-proof pairs per round. However, the lottery evaluation $H(G(s,x), \Pi_i)$ is a single hash---computationally negligible compared to inference and proof generation. Thus, the security-relevant question is whether an adversary can construct \emph{additional valid proof tuples} $\Pi$ from its $q_{\mathsf{ML}}$ inference-proof pairs, thereby extracting extra lottery evaluations without extra inference work. In Bitcoin, this step is trivial: each hash query $H(\mathit{ctr}, G(s,x))$ yields exactly one lottery ticket, and distinct counters give independent evaluations. In PoML, we must rule out that an adversary could obtain additional lottery tickets by reordering proofs, reusing proofs across blocks, forging proofs, or re-randomizing proofs.

\begin{lemma}[Proof-to-Lottery Bijection]\label{lem:bijection}
Let each party complete at most $q_{\mathsf{ML}}$ inference-proof pairs per round, with each proof generated using deterministic prover randomness enforced by the meta-proof $\pi_i^{(2)}$. Then each party can construct at most $q_{\mathsf{ML}}$ distinct valid proof tuples $\Pi$, and consequently obtains at most $q_{\mathsf{ML}}$ lottery evaluations per round. In particular, no adversary can obtain more than $t \cdot q_{\mathsf{ML}}$ lottery evaluations per round across all its controlled parties.
\end{lemma}

\begin{proof}
Fix a miner $(\mathsf{pk}_m, c_{\mathsf{sk}})$ and block header $(s, x)$. We prove by induction on the number of proofs~$k$ the adversary has generated that the set of valid proof tuples has cardinality at most~$k$, except with negligible probability. 

We first establish a key uniqueness property used throughout the induction.

\paragraph{Proof uniqueness.}
For a given inference $M_\theta(c_j, r_j)$ and chain binding $\mathsf{bind}$, there exists at most one valid proof $\pi_j^{(1)}$ with a valid meta-proof $\pi_j^{(2)}$. Indeed, the deterministic randomness $\rho_j = F(\mathsf{sk}_m, r_j)$ fixes $\pi_j^{(1)} = \mathsf{Prove}(\mathsf{CRS},\, \mathsf{stmt}_j^{(1)},\, \mathsf{wit}_j^{(1)}\,;\; \rho_j)$ uniquely. Any alternative $\tilde{\pi} \neq \pi_j^{(1)}$ would require either a valid proof for an incorrect witness (contradicting ZKP soundness), the use of different randomness (contradicting meta-proof soundness), a different block binding (contradicting collision resistance of~$G$), or originating from a different miner (contradicting PRF security and binding of~$c_{\mathsf{sk}}$)---each negligible in $\kappa$.

\paragraph{Base case $(k = 1)$.}
The adversary generates its first proof~$\pi^*$, which has a chain binding $\mathsf{bind}_1 = G(s,x)$ (position 1) baked into its ZKP statement. This creates exactly one valid tuple $(\pi^*)$. By the uniqueness property, no alternative proof $\tilde{\pi} \neq \pi^*$ can form a valid length-1 tuple. Hence there is exactly one valid proof tuple.

\paragraph{Inductive step.}
Assume that after $k$ proofs, the adversary has produced at most $k$ distinct valid proof tuples. Let $\mathcal{S}_k$ denote this set. The adversary generates the $(k+1)$-th proof $\pi^*$. We show that $|\mathcal{S}_{k+1}| \leq k + 1$ by analyzing where $\pi^*$ is bound.

\emph{Case 1: $\pi^*$ is at position $1$} (bound to $G(s,x)$). The new proof forms a length-1 tuple $(\pi^*)$. This is the only new tuple: any longer tuple containing $\pi^*$ as its first element would need a proof at position~2 with binding $H(\pi^*)$. No previously generated proof can have this binding, since producing such a proof before $\pi^*$ exists would require predicting $H(\pi^*)$---a hash preimage, infeasible except with negligible probability. Hence $\pi^*$ contributes exactly one new tuple: $|\mathcal{S}_{k+1}| \leq k + 1$.

\emph{Case 2: $\pi^*$ is at position $i \geq 2$} (bound to $H(\pi_{\mathrm{prev}})$ for some existing proof $\pi_{\mathrm{prev}}$). The binding $\mathsf{bind}_i = H(\pi_{\mathrm{prev}})$ uniquely identifies the predecessor $\pi_{\mathrm{prev}}$ by collision resistance of~$H$. We first argue that all tuples in $\mathcal{S}_k$ containing $\pi_{\mathrm{prev}}$ share the same prefix up to~$\pi_{\mathrm{prev}}$. Indeed, the chain binding of $\pi_{\mathrm{prev}}$ determines its predecessor $\pi_{\mathrm{prev}-1}$, which in turn determines $\pi_{\mathrm{prev}-2}$, and so on back to $G(s,x)$. Therefore, there is a unique chain $\Pi_{\mathrm{prev}} = (\pi_1, \ldots, \pi_{\mathrm{prev}})$ from the root to $\pi_{\mathrm{prev}}$, and the new proof~$\pi^*$ extends it to exactly one new tuple $\Pi_{\mathrm{prev}} \| (\pi^*)$. Since $\pi^*$ is unique for its binding (by the proof uniqueness property), no other extension of this prefix is possible. Hence $|\mathcal{S}_{k+1}| \leq k + 1$.

\paragraph{Conclusion.}
In all cases, each new proof increases the number of valid proof tuples by at most one. After $q_{\mathsf{ML}}$ proofs, there are at most $q_{\mathsf{ML}}$ valid tuples, each yielding one lottery evaluation $H(G(s,x), \Pi)$. Since each corrupted party faces the same constraint, an adversary controlling $t$ parties obtains at most $t \cdot q_{\mathsf{ML}}$ lottery evaluations per round.
\end{proof}

% \begin{remark}[Sequentiality of proof chains]
% The proof chain $(\pi_1, \ldots, \pi_i)$ is sequential: proof $\pi_i$ must be completed before lottery ticket $i$ can be evaluated, since $H_i$ depends on all preceding proofs. Unlike PoW where nonce attempts within a round are parallelizable, a PoML miner cannot ``skip ahead'' to later lottery tickets. However, this sequential structure does \emph{not} expand the adversary's strategy space: the total number of lottery tickets per round remains bounded by $t \cdot q_{\mathsf{ML}}$, exactly as in the backbone's $q$-bounded model. The adversary's only strategic choice is \emph{how to allocate} its $t \cdot q_{\mathsf{ML}}$ proof completions across candidate blocks---a form of selective publishing already captured in the backbone analysis.
% \end{remark}

\subsection{Step 3: The PoML Lottery as a $q_\mathsf{ML}$-Bounded Oracle}

Steps~1 and~2 together show that each flat-model party can produce at most $q_{\mathsf{ML}}$ inference-proof pairs per round (Theorem~\ref{thm:amort-resist}), and that each pair yields exactly one lottery evaluation (Lemma~\ref{lem:bijection}). We now combine these results: define an \emph{inference oracle} as the full pipeline of inference, proof, and hash evaluation on one user inference request. Each party makes at most $q_{\mathsf{ML}}$ queries to this oracle per round. Since the lottery outcome is determined by $H$ (a random oracle), the statistical properties of the PoML lottery are identical to those of Bitcoin's $q$-bounded random oracle.

\begin{lemma}[PoML lottery equivalence]
\label{lem:lottery-equiv}
Under the random oracle model for $H(\cdot)$, the PoML lottery mechanism (Definition~\ref{def:valid-block}) is equivalent to Bitcoin's $q$-bounded random oracle for block validity.
\end{lemma}

\begin{proof}
By Lemma~\ref{lem:bijection}, each party obtains exactly $q_{\mathsf{ML}}$ queries to the inference oracle. Each query evaluates $H_i = H(G(s,x), \Pi)$ and checks $H_i < D$. Since $H$ is modeled as a random oracle with output in $\{0,1\}^\kappa$, each evaluation $H_i$ produces an output uniformly distributed over $\{0,1\}^\kappa$, independent of all previous queries. The per-query success probability is:
\[
    p = \Pr[H(G(s,x), \Pi) < D] = \frac{D}{2^{\kappa}}.
\]
Therefore, each party makes at most $q_{\mathsf{ML}}$ independent random oracle queries per round, each succeeding with probability $p$.

This is structurally identical to the backbone's PoW mechanism where each party makes at most $q$ queries per round to $H$, each succeeding with probability $T / 2^\kappa$. The mapping is:
\[
    q \mapsto q_{\mathsf{ML}}, \qquad T \mapsto D.
\]
\end{proof}


\subsection{Step 4: Deriving $f$ for PoML}

The backbone analysis defines $f$ as the probability that at least one honest party obtains a valid block in a given round. By Lemma~\ref{lem:lottery-equiv}, each of the $(n - t)$ honest parties makes $q_{\mathsf{ML}}$ independent random oracle queries per round, each succeeding with probability $p = D / 2^\kappa$. Therefore:

\[
    f = 1 - (1 - p)^{(n-t) \cdot q_{\mathsf{ML}}}.
\]

For small $p \cdot (n-t) \cdot q_{\mathsf{ML}}$, this approximates to:
\[
    f \approx (n - t) \cdot q_{\mathsf{ML}} \cdot \frac{D}{2^\kappa}.
\]

Since $q_{\mathsf{ML}} \ll q$ (ML inference is orders of magnitude slower than hash computation), the resulting $f$ is naturally smaller for PoML than for PoW at equivalent network sizes. That said, the difficulty parameter $D$ still can be tuned to achieve any desired $f$. For a target block time of $\tau$ seconds with rounds of duration $\Delta$ seconds:
% A smaller $f$ is \emph{favorable} for security: it increases the ratio of uniquely successful rounds ($Y_i = 1$) relative to multiply successful rounds, tightening the bound in Lemma~11 of~\cite{GKL20}. Concretely, the probability that a round is uniquely successful given it is successful approaches~1 as $f \to 0$:
% \[
%     \Pr[Y_i = 1 \mid X_i = 1] = \frac{(n-t) \cdot q_{\mathsf{ML}} \cdot p \cdot (1-p)^{(n-t) \cdot q_{\mathsf{ML}} - 1}}{f}
% \]
% \[
%     \to 1 \quad \text{as } f \to 0.
% \]

%TODO remark 3 of 


\[
    D = \frac{f \cdot 2^\kappa}{(n - t) \cdot q_{\mathsf{ML}}}, \qquad \text{where } f = \frac{\Delta}{\tau}.
\]

% \begin{remark}
% The condition~\eqref{eq:honest-majority} is easier to satisfy in PoML than in PoW precisely because $f$ is smaller. This means PoML can tolerate adversarial fractions closer to $1/2$ for the same network synchronicity assumptions.
% \end{remark}

\subsection{Step 5: Typicality of Executions}

The backbone's security proofs operate within the ``typical execution'' framework (Definition~9, Theorem~10 in~\cite{GKL20}). An execution is \emph{$(\varepsilon, \lambda)$-typical} if, for any set $S$ of at least $\lambda$ consecutive rounds (where $\lambda \geq 2/f$), the following hold:
\begin{enumerate}
    \item[(a)] $(1 - \varepsilon)\mathbb{E}[X(S)] < X(S) < (1 + \varepsilon)\mathbb{E}[X(S)]$ \text{and} \\ $(1 - \varepsilon)\mathbb{E}[Y(S)] < Y(S)$.
    \item[(b)] $Z(S) < \mathbb{E}[Z(S)] + \varepsilon \mathbb{E}[X(S)]$.
    \item[(c)] No insertions, no copies, and no predictions occurred.
\end{enumerate}

Here $X(S) = \sum_{r \in S} X_r$ counts the number of \emph{successful} rounds in $S$ (rounds in which at least one honest party obtains a valid block), $Y(S) = \sum_{r \in S} Y_r$ counts the number of \emph{uniquely successful} rounds (exactly one honest party succeeds), and $Z(S) = \sum_{r \in S} Z_r$ counts the number of rounds in which at least one adversarial party succeeds. In PoML, we adopt the exact same definitions.

\begin{lemma}[Typicality under PoML]
\label{lem:typicality}
PoML executions are $(\varepsilon, \lambda)$-typical (Definition~9 in~\cite{GKL20}) with probability at least
$1 - e^{-\Omega(\varepsilon^2 f \lambda \;+\; \kappa \;-\; \log L)}$.
\end{lemma}

\begin{proof}
By Lemma~\ref{lem:lottery-equiv}, the PoML lottery operates in the $q_{\mathsf{ML}}$-bounded random oracle model with the \emph{same} random oracle $H$ as the backbone protocol. In particular, each lottery query evaluates $H$ on a distinct input (due to the sequentially-bound proof chain), and each evaluation succeeds independently with probability $D/2^\kappa$. This is structurally identical to the backbone's PoW lottery: independent Bernoulli trials under a $q$-bounded random oracle.  Therefore, the random variables $X_r, Y_r, Z_r$ in PoML satisfy the same independence and distributional properties as in the backbone, and the proof of Theorem~10 in~\cite{GKL20} (Chernoff bounds for parts (a)--(b), collision/prediction bounds for part (c)) applies verbatim with $q = q_{\mathsf{ML}}$ and $T = D$.
\end{proof}
% (Note that although PoML evaluations within a round are \emph{sequential}---unlike Bitcoin's parallelizable hash queries---this affects only computation time, not the statistical properties: the Chernoff bounds require independence, not parallelizability, and independence follows from the random oracle model.)

% \begin{remark}
% \red{For the overall probability to be negligible, we need $\varepsilon^2 f \lambda \gg \log L$ . Since $f$ is small (fewer lottery tickets per round), we would require a larger $\lambda$ to compensate. In words: PoML needs longer observation windows before the security guarantees take effect.}
% \end{remark}

% \begin{remark}[Variable inference time]
% In practice, the number of completed inferences per round may vary due to hardware heterogeneity or model complexity. If $q_{\mathsf{ML}}$ is not deterministic but is instead a random variable with bounded support $[q_{\min}, q_{\max}]$, the concentration bounds still hold by conditioning on the worst case $q_{\mathsf{ML}} = q_{\max}$ for the adversary and $q_{\mathsf{ML}} = q_{\min}$ for honest parties. The honest majority condition~\eqref{eq:honest-majority} must then be verified under these worst-case assumptions:
% \[
%     t \cdot q_{\max} < (1 - \gamma) \cdot (n - t) \cdot q_{\min}.
% \]
% This is a strictly stronger requirement than the flat-model assumption, but is necessary to account for hardware variance in practice.
% \end{remark}

\subsection{Main Theorem}

Combining the above results, we obtain the following:

\begin{theorem}[PoML inherits backbone properties]
\label{thm:main-reduction}
Let $\Pi_{\mathsf{PoML}}$ denote the PoML consensus protocol with $n$ parties, at most $t$ of which are corrupted. Each party can complete at most $q_{\mathsf{ML}}$ inference-proof pairs per round. Let $\kappa$ denote the security parameter (hash output length) and $\lambda \geq 2/f$ the tail-bounds parameter. Suppose the model's amortizability parameter satisfies $\delta < 1/(q_{\mathsf{ML}} + 1)$ (Theorem~\ref{thm:amort-resist}). Define:
\[
    f = 1 - (1 - D / 2^\kappa)^{(n-t) \cdot q_{\mathsf{ML}}}.
\]

If the honest majority condition
\[
    t \leq (1 - \gamma)(n - t), \qquad 3f + 3\varepsilon < \gamma,
\]
is satisfied, then $\Pi_{\mathsf{PoML}}$ satisfies the following properties for any $(\varepsilon, \lambda)$-typical execution (which occurs with probability at least $1 - e^{-\Omega(\varepsilon^2 f \lambda \;+\; \kappa \;-\; \log L)}$):
\begin{enumerate}
    \item \textbf{Common Prefix} with parameter $k \geq 2\lambda f$: for any two honest parties with chains $\mathcal{C}_1$, $\mathcal{C}_2$ adopted at rounds $r_1 \leq r_2$, $\mathcal{C}_1^{\lceil k} \preceq \mathcal{C}_2$.
    \item \textbf{Chain Quality} with parameters $\mu$ and $\ell \geq 2\lambda f$: in any $\ell$ consecutive blocks of any honest party's chain, the fraction of honest blocks is at least $\mu$, where $\mu \geq 1 - (1 + \frac{\gamma}{2}) \frac{t}{n - t} \cdot \frac{1}{1 - \gamma}$.
    \item \textbf{Chain Growth} with parameters $\tau$ and $s$: for any $s \geq \lambda$ consecutive rounds, the chain of any honest party grows by at least $\tau \cdot s$ blocks, where $\tau = (1 - \varepsilon)f$.
\end{enumerate}
\end{theorem}

\begin{proof}
By Lemma~\ref{lem:bijection}, each party's $q_{\mathsf{ML}}$ inference-proof completions yield exactly $q_{\mathsf{ML}}$ lottery queries to the random oracle, and the adversary is bounded to $t \cdot q_{\mathsf{ML}}$ queries per round. By Lemma~\ref{lem:lottery-equiv}, each lottery query succeeds independently with probability $D/2^\kappa$, matching the backbone's $q$-bounded random oracle model with $q = q_{\mathsf{ML}}$ and $T = D$. By Lemma~\ref{lem:typicality}, PoML executions are $(\varepsilon, \lambda)$-typical with the stated probability, matching the typicality guarantee of Theorem~10 in~\cite{GKL20}.

Therefore, in any typical execution, $\Pi_{\mathsf{PoML}}$ satisfies the conditions of Lemma~11 in~\cite{GKL20} (with $q_{\mathsf{ML}}$ replacing $q$, $D$ replacing $T$, and $\gamma$ replacing their $\delta$), and Theorems~12 (Chain Growth), 15 (Common Prefix), and~16 (Chain Quality) of~\cite{GKL20} apply directly. The parameter bounds ($k \geq 2\lambda f$, $\ell \geq 2\lambda f$, $\tau = (1 - \varepsilon)f$, and the expression for $\mu$) are taken verbatim from~\cite{GKL20}, with $f$ computed using $q_{\mathsf{ML}}$ and $D$ in place of $q$ and $T$.
\end{proof}


% \subsection{Block Propagation Delay (TODO)}

% PoML blocks contain multiple ZKPs $(\pi_1, \ldots, \pi_k)$ in addition to standard block data. Let $|\pi|$ denote the size of a single proof and $k_{\mathsf{avg}}$ denote the expected number of proofs per block. The total block size is:
% \[
%     |B_{\mathsf{PoML}}| = |B_{\mathsf{header}}| + k_{\mathsf{avg}} \cdot |\pi| + |B_{\mathsf{txns}}|.
% \]

% In the backbone analysis, the round duration must be sufficient for block propagation. Let $\Delta$ denote the maximum network delay (the ``Diffuse'' functionality bound). The critical requirement is that $\Delta$ is small relative to the expected inter-block time $1/f$:
% \[
%     f \cdot \Delta \ll 1.
% \]

% When this product is small, the probability that two honest parties find blocks in the same ``effective round'' (accounting for propagation) is negligible, preserving the uniquely-successful-round analysis.

% \begin{proposition}[Propagation bound for PoML]
% \label{prop:propagation}
% Let $\Delta_{\mathsf{PoW}}$ and $\Delta_{\mathsf{PoML}}$ denote the maximum propagation delays for PoW and PoML blocks respectively, and let $f_{\mathsf{PoW}}$, $f_{\mathsf{PoML}}$ denote the corresponding per-round success probabilities. The backbone security properties hold for PoML provided:
% \[
%     f_{\mathsf{PoML}} \cdot \Delta_{\mathsf{PoML}} \leq f_{\mathsf{PoW}} \cdot \Delta_{\mathsf{PoW}}.
% \]
% \end{proposition}

% \begin{proof}
% The backbone theorems are parameterized by $f$ and implicitly by $\Delta$ (through the round structure). The security guarantees degrade as $f$ increases (cf.\ the requirement $3f + 3\varepsilon < \delta$). Since the effective synchronicity parameter is $f \cdot \Delta$ (the probability of a collision within a propagation window), maintaining $f_{\mathsf{PoML}} \cdot \Delta_{\mathsf{PoML}} \leq f_{\mathsf{PoW}} \cdot \Delta_{\mathsf{PoW}}$ ensures the PoML system is at least as synchronous as the PoW system from the perspective of the backbone analysis.
% \end{proof}

% \begin{remark}[Practical considerations]
% Modern ZKP systems (e.g., SNARK-based provers) produce proofs of size $|\pi| \approx 200$--$500$ bytes. With $k_{\mathsf{avg}} = 10$ proofs per block, the overhead is $\sim$5\,KB, which is negligible compared to typical block sizes ($\sim$1--2\,MB in Bitcoin). Moreover, since $f_{\mathsf{PoML}}$ is naturally small (due to slow inference), the product $f_{\mathsf{PoML}} \cdot \Delta_{\mathsf{PoML}}$ can be made arbitrarily small by tuning $D$. Thus, propagation delay is \textbf{not a practical obstacle} for PoML.
% \end{remark}
